\chapter{状態ベクトル、ユニタリ発展、並進対称性}
\providecommand{\ket}[1]{\left\lvert#1\right\rangle}
\providecommand{\bra}[1]{\left\langle#1\right\rvert}

2四半期の空白の後、量子力学の迅速な復習は有用である。それは四半期全体を占めることはできない。これから先の主題は、量子力学がもともと開発された問題への応用である。すなわち、原子、電子と関連粒子、そして光子である。

原子が最初に来る。電子が原子の一部であるのに、どうして原子が電子より先に来るのか。最初は水素原子の基本方程式を書くのに必要なだけの電子力学があればよい。電子に似た粒子のより詳細な研究は後で行うことができる。これらの応用を通じて、対称性は繰り返し現れる道具となる。すなわち、対称性が量子力学でどのように表現され、系について何を意味するかである。

\section{状態、観測量、直交性}

古典力学では、状態は位相空間の点で表される。量子力学では、状態は複素ベクトル空間のベクトルで表される。これは抽象的なベクトルであり、通常の三次元空間の矢印ではない。状態ベクトルは加算でき、複素数を掛けることができ、それらを含む空間は特定の有限次元を持つ必要はない。

ディラック記法では、状態ベクトルをケットとして書き、
\begin{equation}
\ket{\psi},
\end{equation}
それに対応するブラを関連付ける。
\begin{equation}
\bra{\psi}.
\end{equation}
ブラとケットはエルミート共役によって関係付けられる。ブラとケットを組み合わせると内積が得られ、
\begin{equation}
\langle\phi|\psi\rangle,
\end{equation}
これは一般に複素数である。

観測量とは測定できるものである。位置、運動量、エネルギー、またはその他の測定可能な量である。量子力学では、観測量は線形エルミート演算子\(A\)で表され、
\begin{equation}
A^\dagger=A.
\end{equation}
エルミート性は実数性の演算子版である。特に、エルミート演算子の固有値は実数である。

\(A\)の固有ベクトルは次を満たす。
\begin{equation}
A\ket{\alpha}=\alpha\ket{\alpha}.
\end{equation}
可能な固有値は\(A\)を測定したときの可能な結果である。例えば、観測量が\(1\)、\(3\)、または\(7\)のみを与え得るなら、これらの数がその固有値である。その固有ベクトルは、結果が統計的ではなく確定している状態である。状態\(\ket{\alpha}\)では、\(A\)を測定すると確実に\(\alpha\)が得られる。

内積は状態間の重要な論理関係を表す。\(\ket{\phi}\)と\(\ket{\psi}\)を確実に区別する測定があるとしよう。そのとき状態は直交する。
\begin{equation}
\langle\phi|\psi\rangle=0.
\end{equation}
直交性とは、状態が曖昧さなく物理的に区別可能であることを意味する。規格化された状態では、
\begin{equation}
\langle\psi|\psi\rangle=1.
\end{equation}

\section{位置表示における粒子}

まず、直線上を運動する粒子を考える。その位置\(x\)は観測量である。起こりうる各位置\(x_0\)に対して、位置測定が確実に\(x_0\)を与える状態\(\ket{x_0}\)が存在する。

異なる位置に対応する状態は区別可能であり、したがって直交する。
\begin{equation}
x\neq x'
\quad\Longrightarrow\quad
\langle x'|x\rangle=0.
\end{equation}
任意の状態\(\ket{\psi}\)に対して、その位置状態との重なりが波動関数を定義する。
\begin{equation}
\psi(x)=\langle x|\psi\rangle.
\end{equation}
これが状態ベクトルの位置表示である。

波動関数自体は確率ではない。それは正、負、虚数、または完全に複素数であり得る。位置確率密度は
\begin{equation}
P(x)=\psi^*(x)\psi(x)=|\psi(x)|^2.
\end{equation}
この積は確率密度がそうでなければならないように非負である。規格化された波動関数では、
\begin{equation}
\int_{-\infty}^{\infty}dx\,|\psi(x)|^2=1.
\end{equation}

位置表示では、位置演算子は掛け算として作用する。
\begin{equation}
(\hat x\psi)(x)=x\psi(x).
\end{equation}
その固有関数は一点に集中した理想化された関数である。それらはディラックのデルタ関数として形式的に記述される。すなわち、位置固有値を中心とする非常に狭いスパイクである。それらの正確な超関数的定義は現在の復習には必要ない。

三次元では、位置は三つの成分\(x\)、\(y\)、\(z\)を持ち、位置状態は三次元空間の点を表す。そのような二つの状態は、点が少なくとも一つの座標で異なるときはいつでも直交する。一次元の記法は通常保持され、対応する三次元への拡張は理解されているものとする。

\section{運動量と確定運動量状態}

位置の各成分に対して、対応する運動量の成分が存在する。位置表示では、\(x\)軸に沿った運動量演算子は
\begin{equation}
\hat p_x=-i\hbar\frac{d}{dx}.
\end{equation}
したがって、抽象的なエルミート演算子は波動関数への操作として具体的に見ることもできる。

確定運動量\(p_0\)を持つ状態は固有値方程式を満たす。
\begin{equation}
\hat p_x\psi_{p_0}(x)=p_0\psi_{p_0}(x).
\end{equation}
\(\hat p_x\)の微分形式を代入すると
\begin{equation}
-i\hbar\frac{d\psi_{p_0}(x)}{dx}
=p_0\psi_{p_0}(x).
\end{equation}
波動関数の導関数は波動関数自体に比例するので、解は指数関数の形をとる。
\begin{equation}
\psi_{p_0}(x)\propto
\exp\left(\frac{ip_0x}{\hbar}\right).
\end{equation}

\begin{classroomqa}
\lecturetimestamp{00:24:26}
\audiencequestion{この微分方程式は任意の波動関数に対して成り立つか。}
\lecturerresponse{いいえ。これは確定した運動量をもつ状態に対する固有値方程式である。ほとんどの波動関数はこれを満たさず、この方程式は運動量固有関数を選び出す。}
\end{classroomqa}

位置固有状態と運動量固有状態の区別は明確である。位置固有関数は一点に集中している。運動量固有関数は
\begin{equation}
\left|\exp\left(\frac{ip_0x}{\hbar}\right)\right|^2=1.
\end{equation}
をもつ。したがってその絶対値は直線全体にわたって一定である。この理想化された意味で、確定した運動量は位置の完全な不確定性を意味し、同様に確定した位置は運動量の完全な不確定性を意味する。これは位置--運動量不確定性関係の極端な形である。

\begin{classroomqa}
\lecturetimestamp{00:26:36}
\audiencequestion{運動量が\(p_0\)であると分かっているのに、なぜ依然として確率分布が存在するのか。}
\lecturerresponse{表示されている量は位置測定に関するものであり、別の運動量測定に関するものではない。運動量と位置は異なる観測量である。理想化された運動量固有状態に対して、位置振幅は直線全体にわたって一定の絶対値をもつ。運動量に対する振幅は運動量表示に属し、フーリエ変換によって位置の波動関数と関係づけられる。}
\end{classroomqa}

\begin{classroomqa}
\lecturetimestamp{00:28:02}
\audiencequestion{波動関数がディラックのデルタ関数であるとき、位置の確率はどのように見えるか。}
\lecturerresponse{それは指定された位置に集中している。この理想化されたスパイクは位置の確実性を表している。}
\end{classroomqa}

\editorialnote{厳密な位置固有状態と運動量固有状態は一般化された状態である。ディラックのデルタと平面波は、直線全体上の通常の二乗可積分な規格化された波動関数ではなく、デルタ関数の文字通りの二乗は通常の超関数ではない。上述の記述は、それらの理想化された局在と一定の絶対値に関するものである。}

波動関数は形式的には、\(x\)の各値に対して一つの成分をもつベクトルとみなすことができる。しかし、通常の有限の列ベクトルとは異なり、それは連続無限個の成分をもつ。

\begin{classroomqa}
\lecturetimestamp{00:28:53}
\audiencequestion{状態ベクトルは列ベクトルで表されるのか。}
\lecturerresponse{そのように表すこともできるが、\(x\)が連続的に多くの値をとるため、位置の波動関数を通常の列として書くのは厄介である。形式的には連続的な列ベクトルとみなすことができ、対応するブラは連続的な行ベクトルとみなされる。}
\end{classroomqa}

\section{時間発展とユニタリ性}

状態は時間とともにどのように変化するか。時間発展は状態の空間上の変換である。経過時間\(t\)にわたる発展を\(U(t)\)で表す。すると
\begin{equation}
U(t)\ket{\psi(t_1)}=\ket{\psi(t_1+t)}.
\end{equation}
特に、
\begin{equation}
\ket{\psi(t)}=U(t)\ket{\psi(0)},
\qquad
U(0)=\mathbf 1.
\end{equation}

\(U(t)\)に対する条件は情報の保存から従う。二つの初期状態が区別可能であるならば、時間発展はその区別を消してはならない。直交する状態は直交したままでなければならない。より一般に、任意の二つの状態の間の内積が保存される：
\begin{equation}
\langle U(t)\phi|U(t)\psi\rangle
=\langle\phi|\psi\rangle.
\end{equation}

\begin{equation}
\langle U(t)\phi|U(t)\psi\rangle
=\langle\phi|U^\dagger(t)U(t)|\psi\rangle,
\end{equation}
であるから、すべての内積の保存は
\begin{equation}
U^\dagger(t)U(t)=\mathbf 1.
\end{equation}
を要求する。この方程式を満たす演算子はユニタリである。

\begin{classroomqa}
\lecturetimestamp{00:36:36}
\audiencequestion{すべての内積を保存することは、単に直交性を保存することよりも強い主張か。}
\lecturerresponse{ここで議論している線形な状態変換については、直交性の保存は内積の保存に導く。通常の回転が身近な類推を与える：それらは直交性と内積を保存する。}
\end{classroomqa}

\editorialnote{厳密には、直交性の保存だけでは全体の非零の再スケーリングが許される。物理的な時間発展にすでに要求される規格化の保存がその自由度を除き、内積の保存を与える。}

ユニタリ演算子をエルミート観測量と混同してはならない。位置と運動量はエルミートであるが、一般にユニタリではない。ユニタリ演算子の固有値は位相であり、単位円上にある\(e^{i\theta}\)の形の数である。

\begin{classroomqa}
\lecturetimestamp{00:38:12}
\audiencequestion{非ユニタリ変換の例は何ですか。}
\lecturerresponse{位置演算子と運動量演算子がその例である。これらはエルミートな観測量であるが、ユニタリではない。ユニタリ演算子の固有値は絶対値が1であり、したがって位相因子である。}
\end{classroomqa}


\section{ハミルトニアンとSchr\"odinger方程式}

時間発展は連続であると仮定する。非常に短い区間 \(\epsilon\) の間では、発展演算子はしたがって恒等演算子に近くなければならない。
\begin{equation}
U(\epsilon)=\mathbf 1+\epsilon G+\mathcal O(\epsilon^2).
\end{equation}
ユニタリ性を用いると、
\begin{align}
U^\dagger(\epsilon)U(\epsilon)
&=
\left(\mathbf 1+\epsilon G^\dagger\right)
\left(\mathbf 1+\epsilon G\right)
+\mathcal O(\epsilon^2)
\\
&=
\mathbf 1+\epsilon(G^\dagger+G)
+\mathcal O(\epsilon^2).
\end{align}
一次まででは、
\begin{equation}
G^\dagger=-G.
\end{equation}
したがって \(G\) は反エルミートである。これは
\begin{equation}
G=-\frac{i}{\hbar}H,
\end{equation}
と書ける。ここで \(H\) はエルミートである。
\begin{equation}
H^\dagger=H.
\end{equation}
したがって無限小時間発展演算子は
\begin{equation}
U(\epsilon)
=
\mathbf 1-\frac{i\epsilon}{\hbar}H
+\mathcal O(\epsilon^2).
\end{equation}
である。エルミート演算子 \(H\) がハミルトニアンである。

\begin{classroomqa}
\lecturetimestamp{00:45:48}
\audiencequestion{何が \(\epsilon^2\) の項を無視することを可能にするのですか。}
\lecturerresponse{計算は \(\epsilon\) の一次まで一貫して行われている。二次の計算ならば、\(-\epsilon^2H^2/(2\hbar^2)\) に比例する項を含むすべての二次の項が必要であろう。近似は \(U(\epsilon)\) の打ち切られた表式にあり、\(H\) のエルミート性にあるのではない。}
\end{classroomqa}

有限の時間間隔は多くの小さな区間から構成できる。\(t\)を長さ\(t/N\)の\(N\)個のステップに分割する。無限小変換を繰り返すと
\begin{equation}
U(t)
=
\lim_{N\to\infty}
\left(
\mathbf 1-\frac{itH}{N\hbar}
\right)^N
=
\exp\left(-\frac{iHt}{\hbar}\right).
\end{equation}
が得られる。指数関数は無限小変換の有限版である。

次に状態を\(t\)から\(t+\epsilon\)へ発展させる。
\begin{align}
\ket{\psi(t+\epsilon)}
&=
U(\epsilon)\ket{\psi(t)}
\\
&=
\left(
\mathbf 1-\frac{i\epsilon}{\hbar}H
\right)\ket{\psi(t)}
+\mathcal O(\epsilon^2).
\end{align}
\(\ket{\psi(t)}\)を引き、\(\epsilon\)で割り、極限を取ると
\begin{equation}
\frac{d}{dt}\ket{\psi(t)}
=
-\frac{i}{\hbar}H\ket{\psi(t)}.
\end{equation}
が得られる。等価的に、
\begin{equation}
i\hbar\frac{d}{dt}\ket{\psi(t)}
=
H\ket{\psi(t)}.
\end{equation}
である。これが時間依存Schr\"odinger方程式である。位置表示では
\begin{equation}
i\hbar\frac{\partial}{\partial t}\psi(x,t)
=
H\psi(x,t),
\end{equation}
となる。ここで\(H\)は波動関数に作用する演算子として働く。

時間に依存しないSchr\"odinger方程式はハミルトニアンの固有値方程式である。
\begin{equation}
H\ket{E}=E\ket{E}.
\end{equation}
固有値\(E\)は測定されうるエネルギー、すなわちエネルギー準位である。したがって同じハミルトニアンが、許されるエネルギーと状態の発展の両方を決定する。

初期状態がエネルギー固有状態ならば、その発展は
\begin{equation}
\ket{E;t}
=
\exp\left(-\frac{iEt}{\hbar}\right)\ket{E}.
\end{equation}
である。それは同じエネルギー固有ベクトルのままであり、位相だけを得る。\footnote{負号は\(U(t)=\exp(-iHt/\hbar)\)の先行する導出によって固定される。後に口頭で述べられた位相の表式は反対の符号を持ち、その導出と矛盾する。}

\begin{classroomqa}
\lecturetimestamp{01:01:05}
\audiencequestion{エネルギー固有ベクトルは時間依存Schr\"odinger方程式の下でどのように発展するか。}
\lecturerresponse{ケットを保つ。\(H\ket{E}=E\ket{E}\)であるから、時間発展はエネルギー固有ベクトルにそのエネルギーで決まる位相を掛ける。\(\exp(-iEt/\hbar)\)。}
\end{classroomqa}

\section{何が対称性とみなされるか。}

時間発展は状態空間上の一つのユニタリ変換である。対称変換はもう一つの重要なクラスである。

回転対称性は原子物理学の中心である。水素原子では、電子は中心クーロンポテンシャル中を運動する。座標系を回転しても方程式は変わらず、物理的に原子を回転しても同じ方程式に支配される水素原子のままである。状態は回転の下で変化しうるが、力学法則は変化しない。

並進はもう一つの例を与える。空空間中の孤立系は、そのすべての部分が同じ量だけ変位した後でも同じように振る舞うはずである。この限定条件は重要である。炉が存在し、系だけが炉に近づけられるならば、系は異なる振る舞いをするかもしれない。炉はその並進対称性を破る。系と炉の両方が一緒に並進されるならば、それらの間の元の関係は回復される。

すべての操作が対称性であるわけではない。岩塩の結晶を半分に圧縮すると、一般にその振る舞いは変わる。圧縮は明確に定義された変換であるが、一般には対称性ではない。

したがって対称性とは、力学の形を保つ状態上の操作である。それを\(V\)で表す。それは識別可能性を保たなければならない。2つの直交状態は変換後も直交しなければならない。したがって、
\begin{equation}
V^\dagger V=\mathbf 1,
\end{equation}
であるから、\(V\)はユニタリである。

\section{対称性は時間発展と可換である}

初期状態\(\ket{\psi_1}\)が\(\ket{\psi_2}\)へ発展するとする。
\begin{equation}
\ket{\psi_2}=U(t)\ket{\psi_1}.
\end{equation}
提案された対称性\(V\)を両方の状態に適用する。
\begin{equation}
\ket{\psi_1'}=V\ket{\psi_1},
\qquad
\ket{\psi_2'}=V\ket{\psi_2}.
\end{equation}
\(V\)が力学の対称性ならば、変換された初期状態は変換された最終状態へ発展しなければならない。
\begin{equation}
\ket{\psi_2'}=U(t)\ket{\psi_1'}.
\end{equation}
定義を代入すると
\begin{equation}
VU(t)\ket{\psi_1}
=
U(t)V\ket{\psi_1}.
\end{equation}
が得られる。これはすべての初期状態について成り立たなければならない。したがって
\begin{equation}
VU(t)=U(t)V.
\end{equation}

二つの等価な経路は可換図式として表示できる。\footnote{この図式は01：10：30と01：16：34の間で展開された二経路の議論を再構成する。}
\begin{figure}
\centering
\begin{tikzpicture}[>=Stealth]
\node (a) at (0,0) {$\ket{\psi_1}$};
\node (b) at (4.2,0) {$\ket{\psi_2}=U(t)\ket{\psi_1}$};
\node (c) at (0,-2.0) {$V\ket{\psi_1}$};
\node (d) at (4.2,-2.0) {$V\ket{\psi_2}=U(t)V\ket{\psi_1}$};

\draw[->] (a) -- node[above] {$U(t)$} (b);
\draw[->] (c) -- node[below] {$U(t)$} (d);
\draw[->] (a) -- node[left] {$V$} (c);
\draw[->] (b) -- node[right] {$V$} (d);
\end{tikzpicture}
\caption{対称操作と時間発展は、どちらの順序でも同じ最終状態を与える。}
\end{figure}

\begin{classroomqa}
\lecturetimestamp{01:16:38}
\audiencequestion{対称性はすべての時間について時間発展と可換でなければならないか？}
\lecturerresponse{はい。力学の対称性は、全時間発展に対して可換条件を満たさなければならない。}
\end{classroomqa}

\begin{classroomqa}
\lecturetimestamp{01:17:13}
\audiencequestion{これは、回転してから待ってもよいし、待ってから回転してもよいという意味か？}
\lecturerresponse{そのとおり。回転が対称性ならば、どちらの順序でも同じ変換された最終状態が得られる。}
\end{classroomqa}

小さな時間間隔に対して、
\begin{equation}
U(\epsilon)
=
\mathbf 1-\frac{i\epsilon}{\hbar}H
+\mathcal O(\epsilon^2).
\end{equation}
可換条件は
\begin{equation}
V\left(
\mathbf 1-\frac{i\epsilon}{\hbar}H
\right)
=
\left(
\mathbf 1-\frac{i\epsilon}{\hbar}H
\right)V
+\mathcal O(\epsilon^2).
\end{equation}
となる。恒等項は相殺し、
\begin{equation}
VH=HV,
\end{equation}
または
\begin{equation}
[V,H]=0.
\end{equation}
が残る。したがって、ユニタリ操作が力学の対称性であるのは、それがハミルトニアンと可換であるまさにそのときに限る。

\begin{classroomqa}
\lecturetimestamp{01:19:29}
\audiencequestion{対称操作\(V\)がユニタリであることを要求するものは何か？}
\lecturerresponse{対称性は状態間の論理的区別を保存しなければならない。互いに排他的な直交状態を、もはや区別できない状態に移してはならない。それらの関係を保存することはユニタリ変換を要求する。}
\end{classroomqa}

明示的な時間依存を持たず\(H\)と可換な演算子は保存される。これは対称性と保存則を結びつける。古典力学では、この関係はネーターの定理によって表される。現在の量子力学的な形では、中心となる判定基準はハミルトニアンとの交換子である。

\section{離散対称性と連続対称性}

対称性は離散的または連続的でありうる。鏡映は離散的である。それは左と右を交換し、両者の間に無限小の鏡映の列は存在しない。二つの同一粒子を交換することは別の離散的操作である。印のない両面物体を裏返すことも同種の例を与える。

回転は連続的である。\(90^\circ\)の回転は、同じ軸のまわりの多くの小さな回転から構成できる。空間並進も同じ理由で連続的である。有限の変位は、多くの微小変位から構成できる。

この小さな変換への分割可能性が、連続対称性に特別な形を与える。恒等変換の近くで、
\begin{equation}
V(\epsilon)
=
\mathbf 1-i\epsilon G+\mathcal O(\epsilon^2),
\end{equation}
と書く。ここで\(\epsilon\)は変換をパラメータ化する。ユニタリ性は
\begin{equation}
G^\dagger=G.
\end{equation}
を意味する。エルミート演算子\(G\)は対称性の生成子と呼ばれる。

変換が対称性ならば、
\begin{equation}
[H,V(\epsilon)]=0.
\end{equation}
である。一次まで展開すると
\begin{equation}
[H,G]=0.
\end{equation}
が得られる。したがって、連続対称性はハミルトニアンと可換なエルミート演算子によって生成される。それらの生成子は保存量である。

``生成される''という語は、ここでは具体的な意味を持つ。有限変換は、有限の時間発展が多くの短い時間ステップから構成されるのとちょうど同じように、無限小変換を繰り返し適用することによって構成される。

\section{並進対称性と運動量}

波動関数を距離\(\epsilon\)だけ右へ能動的に平行移動することを考えよう。もとの波動関数が\(x_0\)にピークをもつなら、平行移動後の波動関数は\(x_0+\epsilon\)にピークをもつはずである。したがって座標\(x\)におけるその値は、\(x-\epsilon\)で評価したもとの波動関数である：
\begin{equation}
(T(\epsilon)\psi)(x)=\psi(x-\epsilon).
\end{equation}
微小変位に対しては、
\begin{equation}
\psi(x-\epsilon)
=
\psi(x)-\epsilon\frac{d\psi(x)}{dx}
+\mathcal O(\epsilon^2).
\end{equation}
したがって
\begin{equation}
T(\epsilon)
=
\mathbf 1-\epsilon\frac{d}{dx}
+\mathcal O(\epsilon^2).
\end{equation}
\begin{equation}
\hat p_x=-i\hbar\frac{d}{dx},
\qquad
\frac{d}{dx}=\frac{i}{\hbar}\hat p_x,
\end{equation}
を用いると、無限小平行移動は
\begin{equation}
T(\epsilon)
=
\mathbf 1-\frac{i\epsilon}{\hbar}\hat p_x
+\mathcal O(\epsilon^2).
\end{equation}
となる。対応する有限平行移動は
\begin{equation}
T(\epsilon)
=
\exp\left(-\frac{i\epsilon}{\hbar}\hat p_x\right).
\end{equation}
である。規約\(V(\epsilon)=\mathbf 1-i\epsilon G+\cdots\)のもとでは、次元をもつ平行移動パラメータにより
\begin{equation}
G_{\mathrm{trans}}=\frac{\hat p_x}{\hbar}.
\end{equation}
が得られる。同値な言い方をすれば、指数の中に\(\hbar\)の因子が現れるとき、運動量は空間平行移動の生成子と呼ばれる。

\begin{classroomqa}
\lecturetimestamp{01:34:08}
\audiencequestion{並進の導出で負符号が失われていないか？}
\lecturerresponse{問題は、並進された波動関数に割り当てられる向きである。右への能動的シフトは\(\psi(x+\epsilon)\)ではなく\(\psi(x-\epsilon)\)で表される。その修正により、\(T(\epsilon)=\mathbf 1-i\epsilon\hat p_x/\hbar+\cdots\)。}
\end{classroomqa}

運動量が保存されるかどうかはハミルトニアンに依存する。一次元を運動する自由粒子に対して、
\begin{equation}
H=\frac{\hat p_x^2}{2m}.
\end{equation}
すべての演算子はそれ自身の任意の関数と可換であるから、
\begin{equation}
[H,\hat p_x]=0.
\end{equation}
したがって自由粒子ハミルトニアンは並進の下で不変であり、\(\hat p_x\)は保存される。物理的には、粒子は初期位置をどこに選んでも同じように振る舞う。初期条件を並進させることは、運動全体を並進させるだけである。

三次元では、
\begin{equation}
H
=
\frac{\hat p_x^2+\hat p_y^2+\hat p_z^2}{2m}.
\end{equation}
各運動量成分は\(H\)と可換である：
\begin{equation}
[H,\hat p_x]=[H,\hat p_y]=[H,\hat p_z]=0.
\end{equation}
各空間方向に独立な並進対称性が存在する。運動量は保存量であると同時に、それらの並進に対応する生成子でもある。

次の場合は回転対称性である。異なる軸のまわりの回転は一般に可換でないため、回転はより込み入っている。すなわち、\(x\)回転に続く\(y\)回転は、逆の順序と同じではない。このことが、非可換な対称性の生成子と、回転の群論的構造をもたらす。
